\section{Regular operations and closure under union}

\subsection{The three regular operations}
Let $A$ and $B$ be languages, not necessarily over the same alphabet. If the
alphabets differ, replace both by $\Sigma_A\cup\Sigma_B$.

\begin{definition}[Regular operations]
\[
\begin{aligned}
  A\cup B
    &=\set{w:w\in A\text{ or }w\in B},\\
  A\concat B
    &=\set{xy:x\in A,\, y\in B},\\
  A^*
    &=\set{x_1x_2\cdots x_k:k\geq 0,\, x_i\in A}.
\end{aligned}
\]
By convention $A^0=\set{\eps}$, so $\eps\in A^*$ always, even if
$\eps\notin A$.
\end{definition}

\begin{example}
Let $A=\set{00,11}$ over $\set{0,1}$ and $B=\set{ab,ba}$ over $\set{a,b}$. Then
\[
\begin{aligned}
  A\cup B
    &=\set{00,11,ab,ba}
    &&\text{(alphabet }\set{0,1,a,b}\text{),}\\
  A\concat B
    &=\set{00ab,00ba,11ab,11ba},\\
  A^*
    &=\set{\eps,00,11,0000,0011,1100,1111,\ldots}.
\end{aligned}
\]
\end{example}

\begin{theorem}
Regular languages are closed under union, concatenation, and star.
\end{theorem}
The union case is proved with a product DFA. Concatenation and star need
nondeterminism; those constructions appear in Chapter~\ref{ch:nfa}.

\subsection{Closure under union: product construction}
\begin{theorem}[Union of regular languages]
If $A$ and $B$ are regular, then $A\cup B$ is regular.
\end{theorem}

\begin{proof}
If $A$ and $B$ use different alphabets $\Sigma_1$ and $\Sigma_2$, work over
$\Sigma=\Sigma_1\cup\Sigma_2$. Let
\[
  M_1=(Q_1,\Sigma,\transition_1,q_1,F_1),\qquad
  M_2=(Q_2,\Sigma,\transition_2,q_2,F_2)
\]
recognize $A$ and $B$. The idea is to run $M_1$ and $M_2$
\emph{simultaneously} on the same input.

Define the product DFA
\[
  M=(Q_1\times Q_2,\Sigma,\transition,(q_1,q_2),F),
\]
where
\[
  \transition\bigl((r,s),a\bigr)
    =\bigl(\transition_1(r,a),\transition_2(s,a)\bigr)
\]
and
\[
  F=\set{(r,s):r\in F_1\text{ or }s\in F_2}.
\]

\textbf{Claim.} $M$ accepts $w\in\Sigma^*$ if and only if $w\in A\cup B$.

\textbf{Subclaim.} For every $w\in\Sigma^*$, after reading $w$ the machine
$M_1$ is in state $r$, $M_2$ is in state $s$, if and only if $M$ is in state
$(r,s)$.

\textit{Proof of the subclaim, by induction on $|w|$.}
If $w=\eps$, then $M_1$ is in $q_1$, $M_2$ is in $q_2$, and $M$ is in
$(q_1,q_2)$.

Let $w=w'a$ with $|w'|=n-1$ and last symbol $a=w_n$. By the inductive
hypothesis, $M$ is in $(r,s)$ after $w'$ if and only if $M_1$ is in $r$ and
$M_2$ is in $s$. The next transitions are
\[
  \transition_1(r,a)=r',\qquad
  \transition_2(s,a)=s',\qquad
  \transition\bigl((r,s),a\bigr)=(r',s').
\]
Hence after $w$ the product is in $(r',s')$ if and only if $M_1$ is in $r'$
and $M_2$ is in $s'$.

\textit{Proof of the claim, assuming the subclaim.}
$M$ accepts $w$ if and only if, after $w$, it is in some $(r,s)\in F$, if and
only if $r\in F_1$ or $s\in F_2$, if and only if $M_1$ or $M_2$ accepts $w$,
if and only if $w\in A\cup B$.
\end{proof}

\begin{figure}[ht]
\centering
\begin{tikzpicture}[node distance=16mm and 22mm]
  \node[box] (in) {read $a$};
  \node[box,below left=of in] (m1) {$M_1$: $r\mapsto r'$};
  \node[box,below right=of in] (m2) {$M_2$: $s\mapsto s'$};
  \node[box,below=24mm of in] (p) {product: $(r',s')$};
  \draw[arrow] (in) -- (m1);
  \draw[arrow] (in) -- (m2);
  \draw[arrow] (m1) -- (p);
  \draw[arrow] (m2) -- (p);
\end{tikzpicture}
\caption{The product DFA stores both component states and updates them in lockstep.}
\end{figure}

\begin{exercise}
Let $M$ be a DFA over $\Sigma$ recognizing $A$. Construct a DFA for the
complement $\overline{A}=\Sigma^*\setminus A$ by swapping accept and
non-accept states. (Missing symbols of a larger alphabet should go to a
permanent rejecting sink.)
\end{exercise}

\subsection{Why concatenation is harder}
To recognize $A\concat B$ a deterministic machine would have to know, while
reading, where the $A$-prefix ends and the $B$-suffix begins. That cut-point
cannot in general be determined with finite memory. Therefore:
\begin{itemize}
  \item a single DFA product does not capture concatenation;
  \item we need \emph{nondeterminism}: the machine may guess the cut.
\end{itemize}
